\documentclass[11pt,a4paper]{article}
\usepackage{turkos}

\begin{document}
\phaseheader{6}{Paging and the Higher-Half Kernel}{Turn on the MMU, move the kernel to 0xC0000000, and build the virtual memory layer every later phase stands on}{34}{42}

\ataglance{2--3 weeks}{5}{Phase 5: the PMM and its tests; confident GDB and QEMU monitor use}{Paging enabled, the kernel
linked and running at \code{0xC0000000}, a direct map of the first 512\,MiB of physical memory, a VMM API
(\code{vmm_map}, \code{vmm_unmap}, \code{vmm_virt_to_phys}), pre-allocated kernel page tables, and a page-fault
handler that explains every fault.}

\section{Why this phase}

Paging is the mechanism behind almost everything a modern OS promises. It lets every process have its own private
address space starting at zero (Phase 10), protects the kernel from user programs, shares one physical copy of memory
between processes and copies it only when written (Phase 11), and lets the kernel find frames anywhere in RAM and make
them look contiguous. This phase turns the MMU on and builds the kernel's side of it. You will not have user
processes yet, but every design choice here decides how easy Phases 7--11 will be.

It is also the phase where most hobby kernels stall. When paging is misconfigured, the CPU can't fetch the next
instruction, which faults, and the fault handler can't be fetched either, so the machine resets with no message.
Work in small steps, check each one in the QEMU monitor before taking the next, and commit after every step that
works.

\section{Concepts you need}

\subsection{From address translation to paging}

\OSTEP builds up to paging in stages, and it is worth following them in order. With \textbf{base and bounds}
(chapter 15) the hardware adds a base to every address and checks a limit; it is simple but wastes the memory between
a program's heap and stack. \textbf{Segmentation} (chapter 16, and your Phase 3 GDT) uses several base/bound pairs,
but variable-sized segments fragment physical memory. \textbf{Paging} (chapter 18) cuts both virtual and physical
memory into fixed 4\,KiB pieces and keeps a table per address space that maps each virtual page to a physical frame.
Fixed sizes mean no external fragmentation, and each page can have its own permissions.

\subsection{Two-level paging on x86}

A flat table with an entry for every 4\,KiB page of a 4\,GiB space would take 4\,MiB per process. x86 uses two levels
instead (Figure~\ref{fig:translate}), which \OSTEP chapter 20 explains as a way to avoid storing entries for unused
regions. CR3 holds the physical address of a \textbf{page directory} of 1024 entries. Each present entry (PDE) points
to a \textbf{page table} of 1024 entries (PTEs), and each present PTE points to a 4\,KiB frame. A 32-bit virtual
address is split into a 10-bit directory index, a 10-bit table index and a 12-bit offset. With the PSE extension
(CR4 bit 4), a PDE can instead map a whole 4\,MiB page directly, with no page table.

\begin{figure}[htbp]
\centering
\begin{tikzpicture}[x=0.4cm, y=1cm, every node/.style={font=\sffamily\scriptsize}]
  \draw[draw=tkNavy!70, fill=tkBlue!12] (0,4) rectangle (10,4.7);
  \draw[draw=tkNavy!70, fill=tkGreen!12] (10,4) rectangle (20,4.7);
  \draw[draw=tkNavy!70, fill=tkAmber!14] (20,4) rectangle (32,4.7);
  \node at (5,4.35) {directory index (bits 31--22)};
  \node at (15,4.35) {table index (bits 21--12)};
  \node at (26,4.35) {offset (bits 11--0)};
  \node[font=\sffamily\scriptsize\bfseries, anchor=east] at (-0.3,4.35) {virtual address};
  % structures
  \node[tkhi, minimum width=1.6cm] (cr3) at (-1.5,1.2) {CR3};
  \node[draw=tkNavy!70, fill=tkBlue!6, minimum width=2.2cm, minimum height=2.3cm, align=center] (pd) at (5,1.2) {page\\directory\\(1024 PDEs)};
  \node[draw=tkNavy!70, fill=tkGreen!6, minimum width=2.2cm, minimum height=2.3cm, align=center] (pt) at (15,1.2) {page\\table\\(1024 PTEs)};
  \node[draw=tkNavy!70, fill=tkAmber!8, minimum width=2.2cm, minimum height=2.3cm, align=center] (fr) at (26,1.2) {4\,KiB\\physical\\frame};
  \draw[tkarrow] (cr3) -- (pd.west);
  \draw[tkarrow] (5,4) -- node[right] {selects PDE} (pd.north);
  \draw[tkarrow] (15,4) -- node[right] {selects PTE} (pt.north);
  \draw[tkarrow] (26,4) -- node[right] {byte in frame} (fr.north);
  \draw[tkarrow] (pd.east) -- node[above, align=center, font=\sffamily\tiny] {PDE:\\frame of table} (pt.west);
  \draw[tkarrow] (pt.east) -- node[above, align=center, font=\sffamily\tiny] {PTE:\\frame of page} (fr.west);
\end{tikzpicture}
\caption{How the MMU translates a 32-bit virtual address with 4\,KiB pages. Every pointer in these tables is a \emph{physical} address.}
\label{fig:translate}
\end{figure}

PDEs and PTEs share one format (Figure~\ref{fig:pte}). The top 20 bits are the frame number and the low 12 bits are
flags. The three you set most often are \textbf{P} (present), \textbf{R/W} (writable) and \textbf{U/S} (accessible
from ring 3). For user access, U/S must be set in both the PDE and the PTE. Bits 9--11 are free for the operating
system; Phase 11 uses one of them to mark copy-on-write pages. By default ring 0 may write even to read-only pages;
setting \textbf{CR0.WP} (bit 16) makes the kernel obey R/W too, which copy-on-write will need.

\begin{figure}[htbp]
\centering
\begin{tikzpicture}[x=0.4cm, y=1cm]
  \draw[draw=tkNavy!70, fill=tkGray!10] (0,0) rectangle (20,0.7);
  \node[font=\sffamily\scriptsize] at (10,0.35) {physical frame number (bits 31--12)};
  \draw[draw=tkNavy!70, fill=tkGreen!15] (20,0) rectangle (23,0.7); \node[font=\sffamily\scriptsize] at (21.5,0.35) {AVL};
  \foreach \x/\l in {23/G, 24/PS, 25/D, 26/A, 27/CD, 28/WT} {
    \draw[draw=tkNavy!70, fill=tkGray!10] (\x,0) rectangle (\x+1,0.7);
    \node[font=\sffamily\tiny] at (\x+0.5,0.35) {\l};
  }
  \foreach \x/\l in {29/U, 30/W, 31/P} {
    \draw[draw=tkNavy!70, fill=tkRed!12] (\x,0) rectangle (\x+1,0.7);
    \node[font=\sffamily\scriptsize\bfseries] at (\x+0.5,0.35) {\l};
  }
  \foreach \x/\b in {0.5/31, 19.5/12, 20.5/11, 22.5/9, 23.5/8, 24.5/7, 29.5/2, 30.5/1, 31.5/0}
    \node[font=\sffamily\tiny, text=tkGray] at (\x,0.95) {\b};
  \node[tknote, text width=12.6cm, align=left, anchor=north west] at (0,-0.1)
    {P present \quad W writable \quad U user-accessible \quad WT, CD cache control (PWT, PCD) \quad A accessed \quad D dirty (PTE)
     \quad PS 4\,MiB page (PDE) \quad G global \quad AVL free for the OS};
\end{tikzpicture}
\caption{The format of a page directory or page table entry.}
\label{fig:pte}
\end{figure}

\subsection{The TLB}

Walking two tables in memory for every access would be far too slow, so the CPU caches recent translations in the
\textbf{TLB} (\OSTEP chapter 19). The TLB is \emph{not} kept coherent with your tables automatically. Whenever you
change or remove an existing mapping, you must flush it: \code{invlpg [addr]} for one page, or reload CR3 to flush
everything that is not marked global. Forgetting this is the classic source of heisenbugs: the old mapping keeps
working for a while, and then it doesn't.

\subsection{Why a higher-half kernel}

Once user processes exist, each gets its own page directory, but the kernel must be reachable from all of them:
interrupts and system calls arrive while a user address space is active. The standard solution is to map the kernel
into the top of \emph{every} address space and give the bottom to user programs. Turk-OS uses the Linux-style split:
0--3\,GiB for users and 3--4\,GiB for the kernel (Figure~\ref{fig:layout}). The kernel is linked to run at
\code{0xC0100000} but still loaded by GRUB at physical \code{0x00100000}; \LOB chapter 9 explains why not to simply
identity-map the kernel low instead.

\begin{figure}[htbp]
\centering
\begin{tikzpicture}[m/.style={draw=tkNavy!70, minimum width=7.2cm, font=\sffamily\scriptsize, align=center},
  a/.style={font=\ttfamily\scriptsize, anchor=east}, n/.style={font=\sffamily\scriptsize, anchor=west, text=tkGray, align=left}]
  \node[m, minimum height=0.5cm, fill=tkGray!15] (res) at (0,0) {reserved (recursive mapping, if you ever want one)};
  \node[m, minimum height=0.8cm, fill=tkAmber!12, below=0pt of res] (vm) {kernel stacks and other kernel mappings};
  \node[m, minimum height=0.8cm, fill=tkAmber!12, below=0pt of vm] (heap) {kernel heap (Phase 7)};
  \node[m, minimum height=1.3cm, fill=tkRed!10, below=0pt of heap] (dm) {\textbf{direct map} of physical 0--512\,MiB\\kernel image at \code{0xC0100000}, VGA at \code{0xC00B8000}};
  \node[m, minimum height=2.6cm, fill=tkGreen!8, below=0pt of dm] (user) {user space, private to each process (Phase 10)\\empty for now: address 0 is unmapped, so \code{NULL} faults};
  \node[a] at ([xshift=-4pt]res.north west) {0xFFFFFFFF};
  \node[a] at ([xshift=-4pt]res.south west) {0xFFC00000};
  \node[a] at ([xshift=-4pt]vm.south west) {0xF0000000};
  \node[a] at ([xshift=-4pt]heap.south west) {0xE0000000};
  \node[a] at ([xshift=-4pt]dm.south west) {0xC0000000};
  \node[a] at ([xshift=-4pt]user.south west) {0x00000000};
  \draw[decorate, decoration={brace, amplitude=5pt}, tkRed] ([xshift=4pt]res.north east) -- ([xshift=4pt]dm.south east)
    node[midway, right=6pt, n] {kernel half: the same page\\tables in every process\\(PDEs 768--1023)};
  \draw[decorate, decoration={brace, amplitude=5pt}, tkGreen!70!black] ([xshift=4pt]user.north east) -- ([xshift=4pt]user.south east)
    node[midway, right=6pt, n] {user half\\(PDEs 0--767)};
\end{tikzpicture}
\caption{The Turk-OS virtual address space.}
\label{fig:layout}
\end{figure}

\subsection{Reaching physical memory once paging is on}

After paging is enabled, the kernel can only touch memory that is mapped. Yet it constantly needs to write to frames
it has just taken from the PMM: to zero a new page table, to copy a page, to load a program. \LOB chapter 10 lists the
options. You can \textbf{map frames temporarily} into a reserved slot each time. You can use a \textbf{recursive
mapping}, where the last PDE points at the directory itself, so the tables appear at \code{0xFFC00000}. Or you can
keep a \textbf{direct map} of physical memory in the kernel half, so physical address $p$ is always visible at
$p + \code{0xC0000000}$. Turk-OS uses the direct map, like xv6 and 32-bit Linux: it makes every later phase simpler,
at the price of limiting usable RAM to 512\,MiB, which is plenty for this project. Two macros, \code{P2V} and
\code{V2P}, convert between the two kinds of address.

\section{Reading plan}

\begin{readingplan}
\must & \LOB & Ch.~8 \emph{A Short Introduction to Virtual Memory}; Ch.~9 \emph{Paging} (identity paging, enabling paging, higher-half linker script, entering and running in the higher half) & 51--59 & The practical route for Tasks 6.1--6.4. \\
\must & \OSTEP & Ch.~15 \emph{Mechanism: Address Translation} & 135--148 & What translation is, before you build it. \\
\must & \OSTEP & Ch.~18 \emph{Paging: Introduction}; Ch.~19 \emph{Paging: Faster Translations (TLBs)} & 179--210 & Page tables, PTE bits, why the TLB exists and when to flush it. \\
\must & \OSTEP & Ch.~20 \emph{Paging: Smaller Tables} & 211--226 & Multi-level tables: exactly the x86 design. \\
\should & \OSC & \S9.3 \emph{Paging}; \S9.4 \emph{Structure of the Page Table} & 360--376 & Second explanation with worked examples. \\
\should & \OSC & \S9.6 \emph{Example: Intel 32- and 64-bit Architectures} & 379--383 & The real IA-32 paging hardware, including PAE. \\
\should & \MOS & \S3.3 \emph{Virtual Memory} (paging, page tables, speeding up paging, large memories) & 194--209 & The same ground with more hardware detail. \\
\could & \OSDI & \S4.3 \emph{Virtual Memory} & 383--396 & Another careful walk through page tables and TLBs. \\
\could & \OSTEP & Ch.~23 \emph{The VAX/VMS Virtual Memory System} (now \emph{Complete Virtual Memory Systems}, \S23.1) & 255--264 & A real system that also maps the kernel into every address space. \\
\end{readingplan}

Keep the Intel SDM Volume 3A chapter on paging open, specifically the part on 32-bit paging: it has the exact PDE
and PTE formats, the CR0/CR3/CR4 bits, and the page-fault error code.

\begin{minixbox}
Here MINIX is a contrast rather than a model. The MINIX 3 version in \OSDI does not use paging: each process gets
contiguous text and data segments described by its own local descriptor table, and protection comes from segment
limits (\OSDI\,\S4.6.2 and \S4.7.1). Later MINIX releases added paging, for the reasons \MOS\,\S3.7.3 gives for
why x86 systems abandoned segmentation. If you want a small paging kernel to compare with, look at the MIT xv6
source (\code{vm.c}), which uses the same direct-map approach as Turk-OS.
\end{minixbox}

\section{Build it}

\task{Warm-up: identity paging from C}
Before moving anything, prove you can turn paging on without changing any address. In C, take one frame for a page
directory and one for a page table, map the first 4\,MiB to itself (virtual = physical) with 4\,KiB pages, load CR3,
and set CR0.PG. If the kernel keeps running, your table format is right. You will throw this code away in Task~6.3,
but it isolates the ``is my PTE format correct?'' question from everything else.

\begin{lst}{c}{kernel/arch/i386/paging.c (warm-up only)}
static inline void load_cr3(uintptr_t pd_phys) { __asm__ volatile ("mov %0, %%cr3" : : "r"(pd_phys) : "memory"); }

static inline void enable_paging(void)
{
    uint32_t cr0;
    __asm__ volatile ("mov %%cr0, %0" : "=r"(cr0));
    cr0 |= 0x80000000u;                     /* PG: paging on */
    __asm__ volatile ("mov %0, %%cr0" : : "r"(cr0) : "memory");
}
\end{lst}
\verify{After \code{enable_paging()}, \code{kprintf} still works, and the QEMU monitor's \code{info mem} shows one
range \code{00000000-00400000} marked \code{-rw}.}

\task{A page-fault handler that explains itself}
Register a handler for vector 14 through \code{register_interrupt_handler}. Read CR2 and decode the error code:
bit 0 tells a non-present page from a protection violation, bit 1 a write from a read, bit 2 user mode from kernel
mode, bit 3 a reserved-bit error in a table entry. For now every page fault is fatal: print the decoded message and
panic. From Phase 10 on, faults from user mode will kill only the process, and Phase 11 will fix some faults instead
of reporting them.

\begin{lst}{c}{kernel/mm/fault.c}
static void page_fault_handler(struct regs *r)
{
    uint32_t addr;
    __asm__ volatile ("mov %%cr2, %0" : "=r"(addr));
    kprintf("Page fault at %08x: %s, %s, %s mode%s  (EIP=%08x)\n", addr,
            (r->err_code & 1) ? "protection violation" : "page not present",
            (r->err_code & 2) ? "write" : "read",
            (r->err_code & 4) ? "user" : "kernel",
            (r->err_code & 8) ? ", reserved bit set" : "", r->eip);
    panic("unhandled page fault");
}
\end{lst}
\verify{With identity paging on, reading \code{*(volatile int *)0x00800000} (beyond the mapped 4\,MiB) prints
``page not present, read, kernel mode'' with that address.}

\task{Move the kernel to the higher half}
Now the real change. Link the kernel at \code{0xC0100000} but tell the linker to place each section at a load address
3\,GiB lower, so GRUB still loads it at 1\,MiB. The entry point must be a physical address, because paging is off
when GRUB jumps to it.

\begin{lst}{plain}{link.ld}
ENTRY(_start)
KERNEL_VMA = 0xC0000000;

SECTIONS
{
    . = 0xC0100000;                                     /* virtual: 3 GiB + 1 MiB */
    kernel_start = .;
    .text   ALIGN(4K) : AT(ADDR(.text)   - KERNEL_VMA) { *(.multiboot) *(.text*) }
    .rodata ALIGN(4K) : AT(ADDR(.rodata) - KERNEL_VMA) { *(.rodata*) }
    .data   ALIGN(4K) : AT(ADDR(.data)   - KERNEL_VMA) { *(.data*) }
    .bss    ALIGN(4K) : AT(ADDR(.bss)    - KERNEL_VMA) { *(COMMON) *(.bss*) }
    kernel_end = .;
}
\end{lst}

The loader builds a boot page directory with 4\,MiB pages. It maps the first 4\,MiB twice: at 0 (so the instructions
right after enabling paging can still be fetched) and at \code{0xC0000000}. The rest of the direct map (up to
512\,MiB) goes in at the same time, so every frame the PMM hands out is reachable from the first moment. After jumping
to the higher half, the loader removes the identity mapping.

\begin{lst}{asm}{boot/loader.s (paging part)}
KERNEL_VMA      equ 0xC0000000
KERNEL_PDE      equ KERNEL_VMA >> 22        ; = 768
DIRECT_MAP_PDES equ 128                     ; 128 x 4 MiB = 512 MiB
PDE_4M_RW       equ 0x83                    ; present | writable | 4 MiB page

global _start
_start equ (loader - KERNEL_VMA)            ; physical entry address for GRUB/QEMU

section .data
align 4096
global boot_page_directory
boot_page_directory:
    dd PDE_4M_RW                            ; PDE 0: identity map 0-4 MiB (temporary)
    times (KERNEL_PDE - 1) dd 0             ; PDEs 1-767: user space, empty
%assign i 0
%rep DIRECT_MAP_PDES                        ; PDEs 768-895: 0xC0000000 + p -> p
    dd (i << 22) | PDE_4M_RW
%assign i i+1
%endrep
    times (1024 - KERNEL_PDE - DIRECT_MAP_PDES) dd 0

section .text
loader:                                     ; paging is OFF: we run at physical addresses
    mov  ecx, (boot_page_directory - KERNEL_VMA)
    mov  cr3, ecx                           ; CR3 needs the PHYSICAL address
    mov  ecx, cr4
    or   ecx, 0x00000010                    ; CR4.PSE: allow 4 MiB pages
    mov  cr4, ecx
    mov  ecx, cr0
    or   ecx, 0x80010000                    ; CR0.PG (paging) and CR0.WP (write protect)
    mov  cr0, ecx
    lea  ecx, [higher_half]                 ; an absolute address above 0xC0000000
    jmp  ecx

higher_half:                                ; from here on we run in the higher half
    mov  dword [boot_page_directory], 0     ; drop the identity mapping of 0-4 MiB
    invlpg [0]
    mov  esp, kernel_stack_top              ; a virtual address now
    add  ebx, KERNEL_VMA                    ; Multiboot info pointer: physical -> virtual
    push ebx
    push eax
    call kmain
\end{lst}
\verify{The kernel boots as before. In the QEMU monitor, \code{info mem} shows \code{c0000000-e0000000} and nothing
below it, and \code{info registers} shows EIP above \code{0xc0100000}.}

\task{Fix every address that assumed physical = virtual}
Add \code{P2V} and \code{V2P} to a header, then go through the kernel and fix each place that used a physical
address directly. The VGA buffer is now at \code{P2V(0xB8000)}. The \code{mmap_addr}, \code{mods_addr} and command
line fields of the Multiboot information are physical, so read them through \code{P2V}. The PMM receives
\code{V2P(kernel_start)} and \code{V2P(kernel_end)} when reserving the kernel image, stores its bitmap at a virtual
address, and ignores any RAM above 512\,MiB.

\begin{lst}{c}{include/mm.h}
#define KERNEL_VMA      0xC0000000u
#define DIRECT_MAP_SIZE (512u * 1024 * 1024)

/* Valid only for physical addresses below DIRECT_MAP_SIZE. */
#define P2V(pa) ((void *)((uintptr_t)(pa) + KERNEL_VMA))
#define V2P(va) ((uintptr_t)(va) - KERNEL_VMA)
\end{lst}
\verify{All Phase 5 tests pass again. \code{mem} reports at most 512\,MiB even with \code{-m 1G}.}

\task{The VMM: mapping 4 KiB pages}
Write \code{kernel/mm/vmm.c}. The core is a function that finds, and optionally creates, the PTE for a virtual
address in a given page directory. On top of it, \code{vmm_map} installs a mapping, \code{vmm_unmap} removes one and
flushes the TLB, and \code{vmm_virt_to_phys} walks the tables. Every function takes the page directory as a parameter,
because in Phase 10 you will use the same code on user address spaces.

\begin{lst}{c}{kernel/mm/vmm.c}
#define PG_PRESENT 0x001
#define PG_WRITE   0x002
#define PG_USER    0x004
#define PG_4MB     0x080
typedef uint32_t pde_t, pte_t;

static pte_t *get_pte(pde_t *pd, uintptr_t va, bool create, uint32_t flags)
{
    pde_t *pde = &pd[va >> 22];
    if (!(*pde & PG_PRESENT)) {
        if (!create)
            return NULL;
        uintptr_t pt = pmm_alloc_frame();
        if (!pt)
            return NULL;
        memset(P2V(pt), 0, PAGE_SIZE);                   /* direct map makes this easy */
        *pde = pt | PG_PRESENT | PG_WRITE | (flags & PG_USER);
    }
    ASSERT(!(*pde & PG_4MB));                             /* not inside a 4 MiB page */
    pte_t *table = P2V(*pde & ~0xFFFu);
    return &table[(va >> 12) & 0x3FF];
}

void vmm_map(pde_t *pd, uintptr_t va, uintptr_t pa, uint32_t flags)
{
    pte_t *pte = get_pte(pd, va, true, flags);
    ASSERT(pte != NULL && !(*pte & PG_PRESENT));          /* never remap silently */
    *pte = (pa & ~0xFFFu) | (flags & 0xFFF) | PG_PRESENT;
    __asm__ volatile ("invlpg (%0)" : : "r"(va) : "memory");
}
\end{lst}
\verify{Map a fresh frame at \code{0xF0000000}, write a pattern through it, and read the same pattern through
\code{P2V} of the frame's physical address. After \code{vmm_unmap}, \code{vmm_virt_to_phys} returns 0 and an access
page-faults with ``page not present''.}

\task{Pre-allocate the kernel's page tables}
In Phase 10 every process gets its own page directory, created by copying PDEs 768--1023 from the kernel's. If the
kernel later creates a new page table in its own directory (say, because the heap grew into a new 4\,MiB region),
existing processes would not see it. Avoid the whole problem now: at boot, allocate an empty page table for every PDE
from \code{0xE0000000} to \code{0xFFBFFFFF} (127 tables, about half a megabyte). After that, kernel PDEs never change,
and copying them is always enough.
\verify{\code{info mem} is unchanged (empty tables map nothing), and a loop over PDEs 896--1022 finds all of them
present.}

\task{Test the fault paths deliberately}
Add \code{crash null} (read address 0), \code{crash kwrite} (write to a page you mapped read-only, which only faults
because CR0.WP is set) and \code{crash unmapped} (read \code{0xF8000000}) to the monitor, and add VMM tests to the
suite: map and unmap many pages, check \code{vmm_virt_to_phys} after each, and check that the PMM free count returns to
its starting value.
\verify{Each crash command gives a correctly decoded page-fault report, and \code{make test} passes.}

\section{Testing and debugging}

The QEMU monitor is your best friend in this phase. \code{info mem} lists mapped virtual ranges with their
permissions, \code{info tlb} lists individual page mappings, and \code{info registers} shows CR0, CR2, CR3 and CR4.
Run with \code{-d int,cpu_reset -no-reboot} so a triple fault stops QEMU and shows the last exceptions instead of
rebooting.

\begin{pitfalls}
Triple fault the instant \code{CR0.PG} is set & The next instruction isn't mapped (no identity mapping), or CR3 got a virtual address & Keep PDE 0 until after the jump; load CR3 with \code{boot_page_directory - KERNEL_VMA}. \\
Triple fault right after the jump to \code{higher_half} & The linker script VMA/LMA split is wrong, or PSE wasn't enabled so \code{0x83} entries were read as table pointers & Check \code{i686-elf-readelf -l}: VirtAddr \code{0xC01...}, PhysAddr \code{0x001...}; check CR4 bit 4. \\
GRUB or QEMU refuses the kernel after the move & The entry point is a virtual address & Use the physical \code{_start} symbol as \code{ENTRY}. \\
Page fault inside the page-fault handler, then reset & The handler, its stack or \code{kprintf}'s data touches unmapped memory & Everything the fault path uses must be in the kernel half, which is always mapped. \\
Old data appears after changing a mapping, or a fault after unmapping arrives late or never & Stale TLB entry & \code{invlpg} after every change to an existing entry. \\
Garbage when reading Multiboot data or boot modules & Physical addresses used as pointers & Read them through \code{P2V}; make sure they lie below 512\,MiB. \\
\end{pitfalls}

\section{Milestone}

\begin{milestonebox}[Virtual memory is on]
\begin{checklist}
  \item The kernel runs at \code{0xC0100000}; \code{info mem} shows no mapping below \code{0xC0000000}.
  \item The direct map covers 512\,MiB, and \code{P2V}/\code{V2P} are used everywhere a physical address is needed.
  \item \code{vmm_map}, \code{vmm_unmap} and \code{vmm_virt_to_phys} work on any page directory and flush the TLB.
  \item Kernel page tables for \code{0xE0000000}--\code{0xFFBFFFFF} are pre-allocated.
  \item A null-pointer read, a write to a read-only page and an unmapped access each produce a clear report.
  \item All PMM and VMM tests pass.
\end{checklist}
\end{milestonebox}

\begin{questionbox}
\begin{enumerate}
  \item Translate \code{0xC0101234} by hand: which PDE index, which PTE index, which offset? With the Turk-OS boot directory, which physical address does it reach?
  \item Why must the code that enables paging be identity-mapped at that moment?
  \item What do CR2 and CR3 hold? Why does CR3 need a physical address?
  \item When exactly must you execute \code{invlpg}? When is reloading CR3 better?
  \item Why do both the PDE and the PTE need the U/S bit for user access?
  \item Why put the kernel in the higher half, and why share the same kernel page tables between all processes?
  \item Compare a direct map, a recursive mapping and temporary mappings. What does each cost?
\end{enumerate}
\end{questionbox}

\section{Going further}

Map the kernel image with 4\,KiB pages instead of one 4\,MiB page, so \code{.text} and \code{.rodata} can be read-only;
with CR0.WP set, a stray write into kernel code then faults immediately instead of corrupting it. Mark kernel mappings
\emph{global} (PTE bit 8, enabled by CR4.PGE) so they survive CR3 reloads, which makes the context switches of Phase 8
cheaper. If you want NX protection (no execution from data pages), read about PAE paging, which has three levels and
64-bit entries.

\nextphase{7}{Kernel Heap and Core Data Structures}
\booklegend
\end{document}
